\chapter{Literature Review and Theoretical Background}
\label{chap:litreview}

\section{Core Concepts and Definitions}
\label{sec:concepts}

This section defines the foundational terminology used throughout the thesis. 
Each term is defined once and used consistently in subsequent chapters.

\subsection{Project Management Office (PMO)}
A Project Management Office (PMO) is a centralized organizational unit responsible 
for defining and maintaining project management standards, governance frameworks, 
and reporting structures across an enterprise \cite{pmi2021}. According to the 
Project Management Institute, PMOs can be classified into three types based on 
their degree of control and influence:

\begin{itemize}
    \item \textbf{Supportive PMO} --- provides templates, best practices, and 
    guidance with minimal direct control over projects;
    \item \textbf{Controlling PMO} --- enforces compliance with project management 
    frameworks and requires the use of specific methodologies;
    \item \textbf{Directive PMO} --- takes direct control of projects, with project 
    managers reporting directly to the PMO.
\end{itemize}

In the context of this thesis, the term PMO refers to any formal organizational 
unit in medium or large Latvian enterprises that fulfills at least a controlling 
or directive function in project governance and reporting.

\subsection{Project Management Information System (PMIS)}
A Project Management Information System (PMIS) is a software-based system or 
combination of tools used to support the planning, execution, monitoring, and 
reporting of projects within a PMO environment \cite{laudon2015}. Common PMIS 
tools include Microsoft Project, Jira, SAP, and Power BI. A PMIS may consist 
of a single integrated platform or multiple specialized tools operating in 
parallel, which introduces the challenge of information system integration.

\subsection{Information System Integration}
Information system integration refers to the process of linking two or more 
software systems so that data flows automatically between them without manual 
intervention \cite{hohpe2004}. In PMO environments, integration typically occurs 
between project tracking tools, financial systems, and reporting platforms. The 
degree of integration can range from fully manual data consolidation to fully 
automated, real-time data synchronization via Application Programming Interfaces 
(APIs) or unified platform ecosystems.

\subsection{Integration Maturity Model}
A maturity model is a structured framework that describes the progressive 
development of a capability across a defined set of levels \cite{paulk1993}. 
In the context of information system integration, a maturity model defines 
the stages through which an organization advances from basic, manual integration 
practices toward fully automated, platform-level integration. This thesis adapts 
the concept of maturity modeling to the specific context of PMO information 
systems, proposing a four-level Integration Maturity Model as the primary 
theoretical framework.

\subsection{Business Intelligence and Dashboard Systems}
Business Intelligence (BI) refers to the technologies, processes, and practices 
used to collect, integrate, analyze, and present business data to support 
decision-making \cite{negash2004}. In PMO environments, BI systems typically 
manifest as centralized dashboards that aggregate data from multiple project 
and financial systems into a unified view. The presence or absence of centralized 
BI systems is used in this thesis as an indicator of governance standardization 
maturity.

\subsection{Data Quality}
Data quality refers to the degree to which data is accurate, complete, consistent, 
and timely for its intended purpose \cite{dama2017}. In the context of PMO 
reporting, data quality is directly influenced by the level of system integration: 
manually consolidated data is more prone to errors and delays, while automated 
integration reduces the risk of inconsistencies. This thesis measures data quality 
through perceived data accuracy and timeliness as reported by PMO representatives.

\subsection{Reporting Automation}
Reporting automation refers to the use of software tools and system integrations 
to generate project status reports, dashboards, and portfolio summaries without 
manual data entry or manipulation \cite{pmi2021}. Automated reporting reduces 
the administrative burden on PMO staff and improves the timeliness and consistency 
of management information. In this thesis, reporting automation is operationalized 
as the extent to which PMO reports are generated automatically from integrated 
systems rather than compiled manually.

\subsection{Medium and Large Enterprises in the Latvian Context}
In accordance with the European Commission enterprise size classification, 
a medium enterprise is defined as an organization employing between 50 and 249 
persons, while a large enterprise employs 250 or more persons \cite{ec2003}. 
This thesis targets medium and large Latvian enterprises as they are more likely 
to have established formal PMO structures and sufficiently complex information 
system landscapes to make integration maturity a relevant concern.
